% !PRINTED 476
\weberpage{476}
Die beiden erzeugenden Permutationen $\tau, \omega$ gehören zur ersten Art (Bd. I, \S.\,153), und folglich ist $P_{168}$ ein Theiler der alternirenden \gesperrt{Permutationsgruppe} von sieben Ziffern.

\begin{center}
\S.\,128.

\gesperrt{Gleichungen siebenten Grades mit einer Gruppe 168\textsuperscript{sten} Grades.}
\end{center}

Wir wenden uns jetzt zu der allgemeinen Theorie der speciellen Art von Gleichungen 7\textsuperscript{ten} Grades, deren \gesperrt{Galois}'sche Gruppe sich auf $P_{168}$ reducirt.

Es seien zunächst $\lambda_0, \lambda_1, \lambda_2, \lambda_3, \lambda_4, \lambda_5, \lambda_6$ beliebige Grössen, und

\weq{1}{\tau = (\lambda_0, \lambda_1, \lambda_2, \lambda_3, \lambda_4, \lambda_5, \lambda_6)}

eine cyklische Permutation 7\textsuperscript{ten} Grades. Führen wir noch das Transpositionspaar

\weq{2}{\omega = (\lambda_1, \lambda_2)\,(\lambda_3, \lambda_5)}

ein, so erzeugen diese beiden Permutationen durch ihre Zusammensetzungen und Wiederholungen die ganze Gruppe $P_{168}$.

Setzt man nach \S.\,72, (17) $\chi$ und $\Theta$ daraus zusammen, so erhält man (da, wie schon oben bemerkt, die Zusammensetzung der Permutationen in umgekehrter Reihenfolge, wie bei den Substitutionen, geschehen muss) die folgenden Permutationen der sieben Indices:

\[
\begin{aligned}
\Theta^3 &= \tau^6 \omega \tau \omega = (1, 2, 5, 3)\,(4, 6),
\qquad \Theta = (1, 3, 5, 2)\,(4, 6) \\
\chi &= \omega \tau^3 \Theta \tau^2 = (1, 4, 2)\,(3, 5, 6),
\end{aligned}
\]

Ausdrücke, die sich auch sehr leicht aus \S.\,127, (7) ableiten lassen. Man sieht, dass $\lambda_0$ durch $\omega, \Theta, \chi$ und folglich durch die ganze in $P_{168}$ enthaltene Octa\"edergruppe $P_{24}$ ungeändert bleibt. Es ist leicht, eine Function der sieben Grössen $\lambda$ zu bilden, die zu der Gruppe $P_{168}$ gehört.

Man nehme z. B. das Product $\lambda_0 \lambda_4 \lambda_6$, das aus den durch $\omega$ unberührt bleibenden $\lambda$ besteht, und bilde die Summe der Producte, die sich daraus durch Anwendung der cyklischen Permutation $\tau$ und ihrer Wiederholungen ergeben:

\weq{3}{v = \lambda_0 \lambda_4 \lambda_6 + \lambda_1 \lambda_5 \lambda_0 + \lambda_2 \lambda_6 \lambda_1 + \lambda_3 \lambda_0 \lambda_2 + \lambda_4 \lambda_1 \lambda_3 + \lambda_5 \lambda_2 \lambda_4 + \lambda_6 \lambda_3 \lambda_5.}
